\subsection{Basic Definitions}

\begin{definition}[Spatial zone]
A member $v$ of a finite set $V$ representing a location distinguishable for the target safety property.
\end{definition}

\begin{definition}[Potential movement edge]
An ordered pair $(v_1, v_2) \in V \times V$ representing a possible transition before guards are applied.
\end{definition}

\begin{definition}[Child capability profile]
A finitely represented capability state $q \in Q$.
\end{definition}

\begin{definition}[Caregiver mode]
A finite type $A$ representing the state of the caregiver.
\end{definition}

\begin{definition}[Barrier state]
A finite type $B$ representing the state of a physical barrier.
\end{definition}

\begin{definition}[Edge guard]
A function determining if a transition is possible given the current capabilities, caregiver mode, barrier states, and assumptions $\Gamma$: $V \times V \to Q \to A \to B \to \Gamma \to \text{Prop}$.
\end{definition}

\subsection{Finite Transition Systems}

\begin{definition}[Domestic safety scenario]
A structure $\Sigma$ containing edges, initial states, hazard states, and guards.
\end{definition}

\begin{definition}[Global state]
The complete state of the system $s \in S$, containing zone, capabilities, caregiver mode, and barrier state.
\end{definition}

\begin{definition}[Transition relation]
A relation $s_1 \to_\Sigma s_2$ defining when one state can transition to another.
\end{definition}

\subsection{Reachability and Safety}

\begin{definition}[Reachability]
Reflexive transitive closure of the transition relation.
\end{definition}

\begin{definition}[Hazard predicate]
$s.\text{zone} \in \Sigma.\text{hazard}$.
\end{definition}

\begin{definition}[Assumption-bounded deterministic safety]
Safety holds if every permitted execution from every initial state avoids every hazard.
\end{definition}

\begin{definition}[Invariant deficit]
A reachable hazard state from an initial state.
\end{definition}

\begin{theorem}[Reachability-invariant equivalence]
Safety holds if and only if there is no initial-to-hazard reachable state.
\end{theorem}

\subsection{Barriers and Interventions}

\begin{definition}[Intervention]
A formal modification to the scenario, modelled as a restriction on the transition relation.
\end{definition}

\begin{definition}[Sound intervention]
An intervention $C$ satisfying $\text{Safe}(\Sigma \text{ restricted by } C, \Gamma)$.
\end{definition}

\begin{theorem}[Barrier monotonicity]
When interventions only remove transitions, adding further blocked transitions preserves an already established safety property.
\end{theorem}

\subsection{Intervention Synthesis and Proof Certificates}

\begin{definition}[Intervention cost]
A cost function over interventions.
\end{definition}

\begin{definition}[Optimal intervention]
A minimum-cost sound intervention.
\end{definition}

\begin{definition}[Conservative abstraction]
A relation where every relevant concrete transition is represented by the abstract model.
\end{definition}

\begin{theorem}[Conservative-abstraction safety transfer]
If the abstract model overapproximates the concrete system and is safe, the represented concrete property follows under the abstraction assumptions.
\end{theorem}

\begin{theorem}[Cut-set sufficiency]
If $C$ intersects every initial-to-hazard path, removing $C$ makes hazards unreachable.
\end{theorem}
